\chapter{Phương pháp luận}
\label{chap:methodology}

Chương này trình bày cơ sở lý thuyết và phương pháp cho từng giai đoạn trong pipeline: lý do chọn cảm biến IMU, các kỹ thuật lọc và phân đoạn tín hiệu, chiến lược trích xuất đặc trưng, thuật toán học máy, và hệ thống tạo mã triển khai. Kiến trúc hệ thống tổng thể và các quyết định thiết kế phần mềm được trình bày trong Chương~\ref{chap:design_implementation}.

\section{Dữ liệu cảm biến quán tính cho nhận dạng hoạt động}
\label{sec:imu_data_rationale}

Cảm biến gia tốc kế đo gia tốc tuyến tính trên ba trục (bao gồm thành phần trọng lực), phản ánh biên độ và tính tuần hoàn của chuyển động. Tuy nhiên, khi thiết bị xoay, vectơ trọng trường chiếu lên các trục thay đổi, gây nhập nhằng giữa chuyển động tịnh tiến và thay đổi định hướng. Cảm biến con quay hồi chuyển đo vận tốc góc (°/s), không bị ảnh hưởng bởi trọng trường và nắm bắt chính xác động học xoay---giúp phân biệt các cặp hoạt động khó như leo/xuống cầu thang hay đi bộ/giậm chân tại chỗ. Tổ hợp sáu trục (3 gia tốc kế + 3 con quay hồi chuyển) cung cấp biểu diễn động học đầy đủ \cite{filippeschi2017survey_imu}, và đây là cấu hình chuẩn trong các tập dữ liệu HAR như UCI-HAR \cite{anguita2013har} và PAMAP2 \cite{reiss2012pamap2}.

Framework không áp đặt tần số lấy mẫu cố định, nhưng có một ràng buộc bắt buộc: \textbf{tần số lấy mẫu phải nhất quán giữa giai đoạn huấn luyện và suy luận triển khai}. Các đặc trưng thống kê được tính theo \textit{số mẫu} trong cửa sổ, không theo thời gian tuyệt đối---nếu tần số thay đổi, cùng một cửa sổ 150 mẫu sẽ đại diện cho khoảng thời gian khác nhau, làm lệch toàn bộ phân phối đặc trưng. Để đảm bảo tính nhất quán, framework nhúng tần số lấy mẫu đã chọn dưới dạng hằng số biên dịch vào mã triển khai. Đối với HAR, dải tần phổ hoạt động phổ biến nằm trong 1--20~Hz \cite{whittle2014gait}, nên tần số lấy mẫu 50--100~Hz là phù hợp.

\section{Quy trình tiền xử lý tương tác}

\subsection{Tải lên và trực quan hóa dữ liệu}
Người dùng tải lên các tập dữ liệu CSV chứa các phép đo IMU, tổ chức theo hoạt động (mỗi tệp tương ứng với một lớp). Framework tự động phát hiện tên cột cảm biến, tần số lấy mẫu và hiển thị biểu đồ đồng bộ của tín hiệu thô trên tất cả các trục, cho phép kiểm tra trực quan chất lượng dữ liệu trước khi tiền xử lý (xem giao diện tại Hình~\ref{fig:ui_data_upload}, Chương~\ref{chap:design_implementation}).

\subsection{Các phương pháp lọc và làm sạch tín hiệu}
\label{sec:signal_filtering}

Dữ liệu IMU thô chứa nhiễu cảm biến, nhiễu điện từ và các xung bất thường cần được loại bỏ trước khi trích xuất đặc trưng. Framework cung cấp năm phương pháp lọc, mỗi phương pháp có đặc tính khác nhau về mức độ làm mịn, tính nhân quả và---quan trọng nhất cho triển khai nhúng---\textit{tính tương đồng huấn luyện-triển khai} (training-deployment parity). Người dùng chọn phương pháp phù hợp qua giao diện tương tác; kết quả lọc được trực quan hóa cùng tín hiệu gốc để đánh giá hiệu quả trước khi phân đoạn.

\subsubsection{Loại bỏ ngoại lệ (outlier removal)}
Xác định và loại bỏ các mẫu có giá trị vượt quá ngưỡng $k\sigma$ so với trung bình cục bộ (mặc định $k=3$). Đây là bước làm sạch thô, xử lý các xung nhiễu điện từ hoặc bất thường phần cứng (spike artifacts) tạo ra giá trị nằm ngoài phạm vi vật lý hợp lệ của cảm biến. Phương pháp này không thay đổi hình dạng tín hiệu mà chỉ loại bỏ các điểm dữ liệu cô lập rõ ràng sai, nên không có vấn đề parity---nó chỉ áp dụng một lần trên tập dữ liệu huấn luyện (offline) và không cần triển khai trên thiết bị.

\subsubsection{Bộ lọc Butterworth low-pass}
Bộ lọc IIR Butterworth bậc $n$ với tần số cắt $f_c$ có đáp ứng biên độ cực đại phẳng (maximally flat):
$$|H(f)|^2 = \frac{1}{1 + (f/f_c)^{2n}}$$
Để tránh lệch pha, kỹ thuật lọc hai chiều (forward-backward filtfilt) được sử dụng, cho đáp ứng $|H_{\text{eff}}|^2 = |H|^4$ với pha bằng không. Cấu hình mặc định: bậc 2, tần số cắt 5~Hz.

\textbf{Parity triển khai}: Vì HAR đệm toàn bộ cửa sổ $W$ mẫu trước khi xử lý, thiết bị thực hiện cùng thao tác lọc hai chiều trên bộ đệm. Các hệ số được tính trước tại thời điểm sinh mã và nhúng dưới dạng hằng số.

\subsubsection{Bộ lọc Savitzky-Golay}
Bộ lọc Savitzky-Golay \cite{savitzky1964smoothing} khớp đa thức bậc $d$ lên cửa sổ $2m+1$ mẫu bằng bình phương tối thiểu, cho kết quả tương đương tích chập FIR với các hệ số $c_j$ được tính từ ma trận Vandermonde. So với trung bình động, phương pháp bảo tồn tốt hơn các đỉnh tín hiệu. Cấu hình mặc định: $m=2$ (cửa sổ 5 mẫu), $d=2$.

\textbf{Parity triển khai}: Các hệ số $c_j$ được tính trước và nhúng vào mã C++ dưới dạng mảng hằng số; tích chập áp dụng trên bộ đệm cửa sổ với mở rộng biên bằng lặp giá trị biên.

\subsubsection{Bộ lọc FFT low-pass (miền tần số)}
Lọc trong miền tần số thực hiện ba bước: DFT, triệt tiêu các bin tần số cao ($k > k_c = \lfloor f_c N / f_s \rfloor$), rồi IDFT. Đặc điểm: đáp ứng tần số lý tưởng (brick-wall) và không méo pha. Chi phí bộ nhớ $\mathcal{O}(N)$.

\textbf{Parity triển khai}: Thiết bị thực thi cùng ba bước DFT~$\to$~masking~$\to$~IDFT trên bộ đệm cửa sổ với cùng $k_c$, dùng số dấu phẩy động 32-bit nhất quán.

\subsubsection{Bộ lọc Kalman (constant-velocity)}
\label{sec:kalman_filter}

Bộ lọc Kalman constant-velocity 1D \cite{kalman1960new} áp dụng cho mỗi kênh cảm biến IMU được đề xuất như giải pháp lọc có \textit{khoảng cách parity bằng không}. Bộ lọc Kalman hoạt động hoàn toàn theo chiều thuận (forward-only, causal) trong cả hai môi trường: Python huấn luyện \textit{và} C++ triển khai---loại bỏ lớp lỗi mà tất cả bốn phương pháp trên gặp phải.

\textbf{Mô hình toán học.} Mỗi kênh cảm biến được mô hình hóa bằng trạng thái gồm vị trí (giá trị cảm biến) và vận tốc:
$$\mathbf{x}_k = \begin{bmatrix} p_k \\ v_k \end{bmatrix}, \quad
F = \begin{bmatrix} 1 & \Delta t \\ 0 & 1 \end{bmatrix}, \quad
H = \begin{bmatrix} 1 & 0 \end{bmatrix}$$

Ma trận hiệp phương sai nhiễu quá trình:
$$Q = q \begin{bmatrix} \frac{\Delta t^3}{3} & \frac{\Delta t^2}{2} \\ \frac{\Delta t^2}{2} & \Delta t \end{bmatrix}$$

Nhiễu đo: $R = r$ (scalar). Hai tham số điều chỉnh được: $q$ (process noise, mặc định $10^{-3}$---nhỏ hơn = mượt hơn) và $r$ (measurement noise, mặc định $10^{-1}$---lớn hơn = mượt hơn). Kalman gain $K_k$ tự động thích ứng theo tín hiệu: tăng khi tín hiệu biến đổi nhanh (bám sát), giảm khi tín hiệu ổn định (lọc mạnh hơn).

\subsubsection{So sánh tổng hợp các phương pháp lọc}

Bảng~\ref{tab:filter_parity_comparison} tổng hợp năm phương pháp theo tiêu chí quan trọng nhất cho hệ thống triển khai nhúng:

\begin{table}[H]
    \centering
    \caption{So sánh các phương pháp lọc theo tính tương đồng huấn luyện-triển khai}
    \label{tab:filter_parity_comparison}
    \small
    \begin{tabular}{@{}p{2.8cm}ccp{4.2cm}@{}}
        \toprule
        Phương pháp & Causal & Parity & Ghi chú triển khai \\
        \midrule
        Loại bỏ ngoại lệ & --- & \checkmark & Chỉ offline (không cần trên thiết bị) \\
        Butterworth LP & Không$^\dag$ & \checkmark & Per-window two-pass filtfilt trên bộ đệm cửa sổ \\
        Savitzky-Golay & Không$^\dag$ & \checkmark & FIR convolution với hệ số tiền tính \\
        FFT low-pass & Không$^\dag$ & \checkmark & DFT trực tiếp trên bộ đệm cửa sổ \\
        \textbf{Kalman filter} & \textbf{Có} & \checkmark & Forward-only; hoạt động cả streaming lẫn windowed \\
        \bottomrule
    \end{tabular}
    \vspace{4pt}
    \begin{minipage}{\linewidth}
    \footnotesize
    $^\dag$Non-causal trong streaming; parity đạt được nhờ HAR đệm toàn bộ cửa sổ trước khi trích xuất đặc trưng, cho phép áp dụng bộ lọc hai chiều trên bộ đệm.
    \end{minipage}
\end{table}

Kết luận thiết kế: tất cả các phương pháp lọc đều đạt parity hoàn toàn khi suy luận HAR thực hiện theo lô cửa sổ, vì quy trình đã đệm toàn bộ $W$ mẫu trước khi trích xuất đặc trưng---ràng buộc non-causal chỉ áp dụng cho streaming theo từng mẫu. Kalman filter có lợi thế bổ sung: hoạt động được trong cả chế độ streaming lẫn lô, phù hợp cho hiển thị tín hiệu thời gian thực trong quá trình thu thập. Loại bỏ ngoại lệ không được nhân bản trên thiết bị vì cần thống kê toàn cục từ toàn bộ bản ghi, trong khi mỗi chu kỳ suy luận chỉ có một cửa sổ 1,5~giây.

\subsection{Phân đoạn và sinh cửa sổ huấn luyện}
\label{sec:windowing_generation}

Máy học trên tín hiệu liên tục đòi hỏi chuyển đổi chuỗi thời gian thành các vector có độ dài cố định. Framework thực hiện qua hai bước: người dùng xác định các đoạn tín hiệu sạch trên biểu đồ, và hệ thống tự động sinh cửa sổ huấn luyện bằng giải thuật cửa sổ trượt.

\subsubsection{Giải thuật cửa sổ trượt tự động}

Cho đoạn tín hiệu $N$ mẫu, kích thước cửa sổ $W$ và độ chồng lấp $o$~(\%), bước trượt $S = W \cdot (1 - o/100)$, số cửa sổ sinh được:
\begin{equation}
    n_w = \left\lfloor \frac{N - W}{S} \right\rfloor + 1
    \label{eq:sliding_window}
\end{equation}
Ví dụ: đoạn 5~giây (500 mẫu ở 100~Hz) với $W = 150$, $o = 50\%$ sinh ra 6 cửa sổ huấn luyện. Overlap 50\% tăng gấp đôi số mẫu từ cùng lượng dữ liệu thô, đồng thời giảm hiệu ứng biên---cùng một sự kiện chuyển động xuất hiện ở vị trí tương đối khác nhau trong các cửa sổ kế tiếp, giúp mô hình học đặc trưng bất biến vị trí. Overlap cao hơn ($>$75\%) tạo ra tương quan cao giữa các cửa sổ liên tiếp, tăng nguy cơ rò rỉ thông tin.

\subsubsection{Lựa chọn đoạn tương tác}

Giao diện kéo thả cho phép người dùng đánh dấu các vùng sạch trên biểu đồ tín hiệu, bỏ qua đoạn chuyển tiếp hoặc nhiễu. Nhiều đoạn trên cùng một tệp được hỗ trợ. Sau khi xác nhận, cửa sổ trượt chạy tự động trên từng đoạn theo Phương trình~\ref{eq:sliding_window}. Nhãn hoạt động kế thừa từ tệp nguồn nhưng có thể gán lại thủ công. Thiết kế này tách bạch hai quyết định: \textit{đâu là dữ liệu đáng tin cậy} (người dùng kiểm soát) và \textit{cách biến dữ liệu thành mẫu huấn luyện} (hệ thống tự động hóa), giảm đáng kể công sức thu thập trong khi bảo tồn kiểm soát chất lượng (xem Hình~\ref{fig:ui_preprocessing}, Chương~\ref{chap:design_implementation}).

\section{Trích xuất đặc trưng}
\label{sec:feature_extraction}

\subsection{Tham số cửa sổ và ràng buộc nhất quán}

Trích xuất đặc trưng thực hiện trên các cửa sổ được sinh theo Mục~\ref{sec:windowing_generation}: $W = 150$ mẫu (1,5 giây ở $f_s = 100$~Hz), bước trượt $S = 75$ (50\% chồng lấp). Ràng buộc cốt lõi: ($W$, $S$, $f_s$) phải đồng nhất giữa huấn luyện và suy luận---thay đổi $W$ sau huấn luyện làm lệch phân phối đặc trưng. Framework nhúng $W$ và $f_s$ dưới dạng hằng số trong mã sinh ra.

\subsection{Các loại trích xuất đặc trưng}
\label{sec:feature_modes}

Sáu loại trích xuất đặc trưng được thiết kế để đáp ứng các yêu cầu khác nhau về số chiều, khả năng triển khai nhúng và bất biến hướng. Bảng~\ref{tab:feature_modes} tóm tắt toàn diện.

\begin{table}[ht]
\centering
\caption{Sáu loại trích xuất đặc trưng: tín hiệu nguồn, cấu thành và đặc tính triển khai}
\label{tab:feature_modes}
\small
\resizebox{\textwidth}{!}{%
\begin{tabular}{@{}p{3.2cm}p{3.8cm}p{3.2cm}ccc@{}}
\toprule
\textbf{Loại} & \textbf{Tín hiệu nguồn} & \textbf{Đặc trưng tính} & \textbf{\#} & \textbf{Nhúng} & \textbf{OI\textsuperscript{*}} \\
\midrule
\texttt{oi\_time\_only}\textsuperscript{1}
  & acc\_mag, gyro\_mag, acc\_jerk\_mag
  & 15 thống kê thời gian / tín hiệu
  & 41 & \checkmark & \checkmark \\[3pt]
\texttt{orientation\_invariant}
  & Như trên + DFT(acc\_mag, gyro\_mag)
  & Như trên + 10 phổ / DFT
  & 63 & \checkmark & \checkmark \\[3pt]
\texttt{time\_domain}
  & aX, aY, aZ, gX, gY, gZ
  & 15 thống kê / trục
  & 90 & \checkmark & $\times$ \\[3pt]
\texttt{all}
  & aX, aY, aZ, gX, gY, gZ
  & 15 thời gian + 11 tần số / trục
  & 156 & $\times$ & $\times$ \\[3pt]
\texttt{frequency\_domain}
  & aX, aY, aZ, gX, gY, gZ
  & 11 đặc trưng phổ / trục
  & 66 & $\times$ & $\times$ \\[3pt]
\texttt{raw}
  & aX, aY, aZ, gX, gY, gZ
  & Giá trị cảm biến thô
  & 6 & \checkmark & $\times$ \\
\bottomrule
\end{tabular}%
}
\vspace{4pt}
\begin{minipage}{\linewidth}
\footnotesize
\textsuperscript{1}\texttt{oi\_time\_only} = \texttt{orientation\_invariant\_time\_only}.\\
\textsuperscript{*}OI = Orientation-Invariant (bất biến hướng).
Tín hiệu tổng hợp: $\text{acc\_mag} = \sqrt{a_X^2+a_Y^2+a_Z^2}$;
$\text{gyro\_mag} = \sqrt{g_X^2+g_Y^2+g_Z^2}$;
$\text{acc\_jerk\_mag} = \tfrac{d}{dt}\,\text{acc\_mag}$.
\end{minipage}
\end{table}

Hai loại \textbf{bất biến hướng} tính đặc trưng trên biên độ tổng hợp (vector magnitude)---đại lượng bất biến phép xoay---phù hợp khi thiết bị đeo không cố định. Loại \texttt{orientation\_invariant} (63 đặc trưng) là khuyến nghị mặc định: cân bằng giữa tính phong phú và triển khai nhúng, vì bộ sinh mã tạo DFT bằng tổng sin/cos trực tiếp không cần thư viện FFT.

Loại \texttt{time\_domain} (90 đặc trưng) áp dụng 15 thống kê lên từng trục riêng lẻ---mỗi đặc trưng là biểu thức dạng đóng xác định, không phụ thuộc trạng thái, đảm bảo bit-identical trên mọi nền tảng:
\begin{itemize}
    \item \textbf{Xu hướng trung tâm}: mean, median
    \item \textbf{Phân tán}: std, range, IQR
    \item \textbf{Cực trị}: min, max, Q1, Q3
    \item \textbf{Hình dạng}: skewness $\gamma = E[(X-\mu)^3]/\sigma^3$, kurtosis $\kappa = E[(X-\mu)^4]/\sigma^4 - 3$
    \item \textbf{Năng lượng}: RMS, signal energy
    \item \textbf{Xấp xỉ tần số}: zero-crossing rate, mean-crossing rate
\end{itemize}

Loại \texttt{all} và \texttt{frequency\_domain} bổ sung đặc trưng phổ mỗi trục nhưng \textbf{không hỗ trợ triển khai nhúng} vì nhạy với sai lệch số, tốn tài nguyên tính toán, và đặc trưng miền thời gian đã đủ hiệu quả trên các tập HAR chuẩn~\cite{banos2014wisdm}. Loại \texttt{raw} truyền trực tiếp giá trị cảm biến, dùng làm đầu vào CNN.

\subsection{Chiến lược đệm cửa sổ}
\label{sec:window_padding}

Các cửa sổ có thể ngắn hơn kích thước mục tiêu khi người dùng chọn đoạn ngắn hoặc cửa sổ cuối không đủ dài.

\textbf{Vấn đề với zero-padding}: Đặt giá trị cảm biến bằng 0 tạo ra magnitude gia tốc $||\mathbf{a}|| = 0$---vật lý bất khả thi vì gia tốc trọng trường luôn $\approx 9{,}81$ m/s². Điều này gây nhiễu nghiêm trọng cho phân phối đặc trưng (min, median, Q1 bị kéo về 0).

\textbf{Giải pháp}: Lặp giá trị cuối cùng đến kích thước mục tiêu: $\mathbf{W}' = [w_1, \ldots, w_k, \underbrace{w_k, \ldots, w_k}_{N-k}]$. Phương pháp bảo toàn tính thống kê tín hiệu và nhất quán với bộ đệm trên thiết bị biên (luôn chứa dữ liệu thực). Cửa sổ có ít hơn 30\% mẫu mục tiêu bị loại bỏ hoàn toàn.

\textit{Phát hiện thực nghiệm}: Zero-padding khiến mô hình luôn dự đoán sai lớp \textit{walking\_downstairs} bất kể hoạt động thực tế, do đặc trưng nằm ngoài phân phối huấn luyện (xem Mục~\ref{sec:training_deployment_parity}).

\subsection{Tăng cường dữ liệu cho huấn luyện}
\label{sec:data_augmentation}

Để cải thiện tính tổng quát hóa và giảm overfitting---đặc biệt với tập dữ liệu nhỏ từ thu thập đối tượng đơn---năm phương pháp tăng cường dữ liệu được áp dụng ở cấp cửa sổ (trước trích xuất đặc trưng) \cite{um2017data, iwana2021empirical}:

\begin{enumerate}
    \item \textbf{Jitter}: Nhiễu Gaussian $\mathcal{N}(0, \sigma^2)$ ($\sigma=0{,}05$), mô phỏng nhiễu cảm biến \cite{um2017data}
    \item \textbf{Scaling}: Nhân hệ số $s \sim \mathcal{N}(1, \sigma_s^2)$, mô phỏng biên độ chuyển động khác nhau
    \item \textbf{Rotation}: Ma trận xoay ngẫu nhiên ($\leq 15^{\circ}$), mô phỏng thay đổi vị trí cảm biến \cite{eyobu2018feature}
    \item \textbf{Time Warping}: Biến đổi phi tuyến trục thời gian qua nội suy spline, mô phỏng tốc độ thực hiện khác nhau \cite{le_guennec2016augmentation}
    \item \textbf{Permutation}: Hoán vị $k=4$ đoạn con, giúp mô hình học đặc trưng cục bộ thay vì phụ thuộc thứ tự toàn cục
\end{enumerate}

\subsubsection{Tăng cường nhận biết lớp}
\label{sec:class_aware_augmentation}

Áp dụng biến đổi mạnh lên hoạt động tĩnh (still, standing, sitting) tạo dao động nhỏ giống hoạt động cường độ thấp, gây nhầm lẫn lớp. Framework giải quyết bằng chiến lược phân biệt:
\begin{itemize}
    \item \textbf{Hoạt động động}: Áp dụng đầy đủ cả năm phương pháp ($\sigma_{\text{jitter}} = 0{,}05$)
    \item \textbf{Hoạt động tĩnh}: Chỉ micro-jitter ($\sigma = 0{,}01$)---đủ tạo đa dạng mà không phá vỡ đặc tính ``gần bất động''
\end{itemize}
Phân loại tĩnh/động dựa trên ngưỡng biên độ dao động và có thể do người dùng chỉ định thủ công. Chiến lược tuân thủ nguyên tắc rằng tăng cường phải bảo toàn các đặc tính phân biệt của mỗi lớp \cite{buda2018systematic}.

\subsection{Ngưỡng tin cậy cho từ chối hoạt động không xác định}
\label{sec:confidence_threshold}

Trong triển khai thực tế, thiết bị có thể ghi nhận chuyển động không thuộc lớp nào đã huấn luyện. Ngưỡng tin cậy $\tau$ (mặc định 0,6) được tích hợp vào mã suy luận: nếu $\max_c P(y=c|\mathbf{x}) < \tau$, dự đoán trả về ``unknown'' thay vì ép buộc một nhãn:
\begin{equation}
    \hat{y} = \begin{cases}
        \arg\max_c P(y=c|\mathbf{x}) & \text{nếu } \max_c P(y=c|\mathbf{x}) \geq \tau \\
        \text{unknown} & \text{ngược lại}
    \end{cases}
\end{equation}

Xác suất từng lớp được tính theo cơ chế có sẵn: softmax cho NN/CNN, tỷ lệ bỏ phiếu cho RF, softmax trên điểm quyết định OvR cho SVM. Giá trị $\tau = 0{,}6$ yêu cầu tin cậy gấp 3 lần xác suất ngẫu nhiên (bài toán 5 lớp: 20\%). So với Edge Impulse (cần huấn luyện khối anomaly detection riêng), cách tiếp cận này không tốn thêm bộ nhớ hay dữ liệu bổ sung.

\section{Huấn luyện mô hình học máy}

\subsection{Các thuật toán học máy áp dụng}
Bốn nhóm thuật toán học máy dưới đây đã được chứng minh hiệu quả trong nhận dạng hoạt động và phù hợp cho triển khai biên, với các đánh đổi khác nhau giữa độ chính xác, yêu cầu dữ liệu và chi phí tài nguyên:

\paragraph{Random Forest (RF)}
Phương pháp ensemble kết hợp nhiều cây quyết định với bagging và ngẫu nhiên hóa đặc trưng \cite{breiman2001random}. Cấu hình điển hình cho bài toán HAR:
\begin{itemize}
    \item Số cây: 50--200 (tham chiếu: 100 cây), độ sâu tối đa 10--30
    \item Cân bằng lớp: \texttt{class\_weight='balanced'} tự động trọng số mẫu tỷ lệ nghịch với tần suất lớp
    \item Tiêu chí phân nhánh: Độ bất thuần Gini hoặc entropy thông tin
\end{itemize}
RF đặc biệt phù hợp cho HAR vì không yêu cầu chuẩn hóa đặc trưng; chịu đựng tốt với đặc trưng không liên quan (mỗi lần phân nhánh chỉ xét $\sqrt{D}$ đặc trưng ngẫu nhiên); kết quả giải thích được qua độ quan trọng đặc trưng; ổn định với nhiễu nhờ cơ chế bỏ phiếu đa số.

\paragraph{Support Vector Machine (SVM)}
Bộ phân loại phân biệt tìm siêu phẳng tối đa hóa biên phân tách giữa các lớp. Cấu hình điển hình:
\begin{itemize}
    \item Kernel: Radial Basis Function (RBF), phù hợp cho dữ liệu phi tuyến; tham số gamma kiểm soát bán kính ảnh hưởng của từng điểm huấn luyện
    \item Chiến lược đa lớp: One-vs-One (OvO) hoặc One-vs-Rest (OvR)
    \item Cân bằng lớp: \texttt{class\_weight='balanced'}
    \item Chính quy hóa: Tham số $C$ kiểm soát đánh đổi biên rộng---vi phạm biên ($C \in \{0{,}1;\ 1;\ 10\}$)
\end{itemize}
SVM phù hợp khi số chiều đặc trưng cao so với số mẫu---điển hình trong HAR với hàng chục đến hàng trăm đặc trưng và vài trăm mẫu. Yêu cầu chuẩn hóa đặc trưng (StandardScaler) trước huấn luyện để đảm bảo các chiều đóng góp đồng đều vào hàm kernel.

\paragraph{Neural Network (NN)}
Perceptron đa lớp (MLP) với lan truyền ngược, phù hợp học các ánh xạ phi tuyến phức tạp. Một kiến trúc tham chiếu cho bài toán 5 lớp với 63 đặc trưng đầu vào:
\begin{itemize}
    \item Kiến trúc: [63 đầu vào] $\rightarrow$ [100 ẩn, ReLU] $\rightarrow$ [num\_classes, Softmax]
    \item Trình tối ưu: Adam ($\eta=0{,}001$, $\beta_1=0{,}9$, $\beta_2=0{,}999$)
    \item Chính quy hóa: Phạt L2 $\alpha=0{,}0001$; dừng sớm (patience=10)
    \item Lịch trình: Tối đa 500 epoch, batch đầy đủ
\end{itemize}
Kiến trúc một lớp ẩn vừa đủ năng lực học mẫu phi tuyến vừa nhỏ gọn cho triển khai nhúng (ma trận trọng số cố định, không phân bổ bộ nhớ động). Kiến trúc tùy chỉnh với nhiều lớp ẩn hơn có thể xây dựng qua PyTorch, với cơ chế xuất trọng số sang định dạng tương thích để tạo mã triển khai.

\paragraph{Convolutional Neural Network (CNN)}
Mạng tích chập 1D xử lý trực tiếp cửa sổ cảm biến thô (end-to-end), học biểu diễn phân cấp từ tín hiệu mà không cần trích xuất đặc trưng thủ công. Một kiến trúc tham chiếu 1D-CNN:
\begin{itemize}
    \item Đầu vào: Cửa sổ thô $W \times C$ ($W=150$ mẫu, $C=6$ kênh cảm biến)
    \item Khối tích chập: 2--3 lớp Conv1D (kernel $3$--$5$, bộ lọc $32$--$64$) + BatchNorm + MaxPool
    \item Phân loại: Flatten + Dense + Dropout ($0{,}3$--$0{,}5$) + Softmax
    \item Trình tối ưu: Adam với OneCycleLR scheduler
    \item Triển khai: Xuất trọng số thành mảng C/C++ hoặc chuyển đổi sang TFLite Micro (INT8)
\end{itemize}
CNN phù hợp khi tập dữ liệu đủ lớn và các đặc trưng thống kê thủ công chưa đủ phân biệt---mô hình học trực tiếp các biểu diễn phân cấp từ tín hiệu thô. Đổi lại, CNN yêu cầu nhiều dữ liệu huấn luyện hơn và bộ nhớ triển khai lớn hơn (500~KB--2~MB so với 50--200~KB cho mô hình cổ điển).

\subsection{Giao thức huấn luyện và đánh giá}

\subsubsection{Chiến lược chia dữ liệu}
Giao thức phân chia ba tập được áp dụng \cite{goodfellow2016deep}: tập huấn luyện (60\%) để khớp tham số, tập xác thực (20\%) để điều chỉnh siêu tham số và lựa chọn mô hình, tập kiểm tra (20\%) chỉ đánh giá một lần duy nhất sau khi huấn luyện hoàn tất. Tất cả các phân chia sử dụng lấy mẫu phân tầng để bảo tồn tỷ lệ lớp. Xác thực chéo 5-fold có thể thực hiện trên tập train+val kết hợp khi cần so sánh thuật toán với dữ liệu hạn chế, trong khi vẫn giữ tập kiểm tra hoàn toàn độc lập.

\subsubsection{Xử lý mất cân bằng lớp}
Tham số \texttt{class\_weight='balanced'} được bật mặc định cho RF và SVM, gán trọng số $w_c \propto 1/n_c$ tỷ lệ nghịch với kích thước lớp. Phân chia phân tầng đảm bảo tỷ lệ lớp nhất quán giữa các tập. Tổ hợp này đảm bảo tất cả các lớp được nhận dạng đáng tin cậy, không chỉ các lớp xuất hiện nhiều nhất.

\subsection{Tối ưu hóa siêu tham số}
Điều chỉnh siêu tham số sử dụng GridSearchCV với xác thực chéo 5-fold. Các không gian tìm kiếm điển hình: RF tìm trên \texttt{n\_estimators} [50--200], \texttt{max\_depth} [10--30]; SVM tìm trên \texttt{C} [0,1--100] và \texttt{gamma} [scale/auto/0,001--0,1]; NN tìm trên kích thước lớp ẩn và tốc độ học. Tối ưu hóa thường cải thiện hiệu suất cơ sở 2--5\% (xem giao diện cấu hình huấn luyện tại Hình~\ref{fig:ui_training}, Chương~\ref{chap:design_implementation}).

\section{Tạo mã cho triển khai biên}

\subsection{Kiến trúc}
Hệ thống con tạo mã dịch các mô hình đã huấn luyện thành các gói triển khai phù hợp với nhiều họ thiết bị và môi trường thực thi khác nhau, thay vì nhắm đến một bo mạch đơn lẻ. Kiến trúc tuân theo mẫu thiết kế Factory, trong đó generator được chọn theo ba trục: \textit{loại mô hình}, \textit{họ nền tảng đích} và \textit{cách tiếp cận triển khai}. Danh sách khóa nền tảng và backend cụ thể được tóm tắt trong Bảng~\ref{tab:deployment_target_matrix}.

Kiến trúc gồm ba lớp chính:
\begin{itemize}
    \item \textbf{Lớp cơ sở (Base Generator)}: cung cấp phần dùng chung gồm sinh tệp, trích xuất đặc trưng, lưu trữ tham số, kiểm tra tương đồng và ví dụ tích hợp.
    \item \textbf{Lớp đặc thù thuật toán}: mỗi thuật toán (RF, SVM, NN, CNN) có generator riêng hiện thực logic suy luận tương ứng.
    \item \textbf{Lớp đặc thù nền tảng}: các generator cho ARM Cortex-M, MicroPython, Zephyr, cùng các backend TFLite Micro và ONNX Runtime, chịu trách nhiệm điều chỉnh mã theo môi trường đích.
\end{itemize}

\begin{table}[H]
    \centering
    \caption{Ma trận đích triển khai được hỗ trợ}
    \label{tab:deployment_target_matrix}
    \small
    \begin{tabular}{@{}p{2.6cm}p{2.8cm}p{2.6cm}p{4.6cm}@{}}
        \toprule
        Nhóm đích & Khóa nền tảng/backend & Gói sinh ra & Mục đích sử dụng chính \\
        \midrule
        Arduino-compatible & arduino, seeed\_xiao, esp32, m5stack, teensy & .h + .cpp + .ino & Quy trình Arduino IDE/PlatformIO cho các bo vi điều khiển phổ biến \\
        ARM Cortex-M thuần & arm\_cortex\_m & .c & Tích hợp bare-metal hoặc CMSIS trên các vi điều khiển ARM \\
        SDK/RTOS native & generic\_c, generic\_cpp, esp\_idf, zephyr & .h + .c/.cpp + ví dụ & Ứng dụng nhúng native, SDK hãng và RTOS \\
        MicroPython & micropython & module .py + ví dụ .py & Nguyên mẫu nhanh hoặc giáo dục trên board hỗ trợ MicroPython \\
        Backend runtime thay thế & tflite\_micro, onnx\_runtime & mã glue + model nhị phân & Mở rộng cho các pipeline cần runtime engine chuẩn hóa \\
        \bottomrule
    \end{tabular}
\end{table}

Chi tiết hiện thực và giao diện tạo mã được trình bày trong Chương~\ref{chap:design_implementation} (Hình~\ref{fig:ui_code_generation}).

\subsection{Lý do chọn tạo mã trực tiếp thay vì Runtime Engine}
\label{sec:codegen_rationale}

Việc lựa chọn phương pháp triển khai mô hình ML lên vi điều khiển có ảnh hưởng trực tiếp đến mức sử dụng tài nguyên, tính minh bạch và phạm vi thuật toán được hỗ trợ. Phần này phân tích ba cách tiếp cận chính và luận giải cho quyết định thiết kế trong bối cảnh HAR biên.

\subsubsection{Các phương pháp triển khai mô hình trên vi điều khiển}

Có ba cách tiếp cận chính để triển khai mô hình ML trên vi điều khiển, mỗi cách thể hiện đánh đổi khác nhau giữa tính tiện lợi, hiệu suất và kiểm soát:

\paragraph{Phương pháp 1: Runtime Engine (TensorFlow Lite Micro, ONNX Micro Runtime)}
Mô hình được chuyển đổi sang định dạng trung gian (\texttt{.tflite}, \texttt{.onnx}), sau đó được tải bởi một runtime engine chung trên thiết bị. Runtime đọc biểu đồ tính toán và thực thi từng toán tử (operator) theo thứ tự.

\textit{Ưu điểm:} Hỗ trợ nhiều kiến trúc mô hình (CNN, LSTM, Transformer); hệ sinh thái lớn với cộng đồng hỗ trợ mạnh; các tối ưu hóa lượng tử hóa tích hợp (INT8, FP16).

\textit{Nhược điểm:} Runtime chiếm 100--300KB flash bổ sung \cite{warden2019tinyml}---đáng kể trên vi điều khiển có 256KB--1MB flash. Interpreter overhead thêm 15--40\% thời gian suy luận so với mã gốc \cite{lin2022mcunet}. Quan trọng nhất, runtime hoạt động như hộp đen: người dùng không thể kiểm tra hoặc xác minh quá trình trích xuất đặc trưng và tính toán bên trong.

\paragraph{Phương pháp 2: Trình biên dịch mô hình (Edge Impulse EON Compiler, Apache TVM, microTVM)}
Mô hình được biên dịch trước (ahead-of-time compilation) sang mã C tối ưu hóa, loại bỏ runtime nhưng vẫn hoạt động qua chuỗi công cụ đóng của nhà cung cấp.

\textit{Ưu điểm:} Không có interpreter overhead; mã tối ưu hóa cho vi kiến trúc cụ thể; Edge Impulse EON Compiler giảm flash 35--55\% so với TFLite Micro (theo số liệu nhà cung cấp \cite{web:Edge_Impulse}).

\textit{Nhược điểm:} Chuỗi công cụ đóng---mã đã tạo bởi EON Compiler khó kiểm toán và không thể sửa đổi. Apache TVM/microTVM yêu cầu kiến thức biên dịch cấp chuyên gia. Hầu hết các trình biên dịch mô hình bắt buộc mô hình đầu vào theo định dạng ONNX hoặc TFLite, loại trừ các mô hình scikit-learn cổ điển (RF, SVM) mà không qua bước chuyển đổi trung gian phức tạp.

\paragraph{Phương pháp 3: Tạo mã trực tiếp (Direct Code Generation)}
Phương pháp này trích xuất trực tiếp các tham số đã huấn luyện từ đối tượng scikit-learn/PyTorch---trọng số và bias của mạng nơ-ron, support vectors và hệ số kernel của SVM, cấu trúc cây nhị phân của Random Forest, trọng số convolution và batch normalization của CNN---rồi tạo mã nguồn (C/C++ hoặc MicroPython) triển khai thuật toán suy luận tương ứng.

\textit{Ưu điểm:} Không phụ thuộc bên ngoài---mã tự chứa hoàn toàn; mức sử dụng bộ nhớ tối thiểu (chỉ chứa trọng số + logic suy luận); minh bạch hoàn toàn---mọi phép tính có thể kiểm tra và xác minh; hỗ trợ trực tiếp các thuật toán ML cổ điển mà không cần chuyển đổi định dạng; xuất được nhiều ngôn ngữ đích (C, C++, MicroPython).

\textit{Nhược điểm:} Yêu cầu triển khai thủ công logic suy luận cho mỗi thuật toán; không hưởng các tối ưu hóa toán tử tự động cấp vi kiến trúc; khó mở rộng sang các kiến trúc phức tạp (CNN nhiều lớp, attention).

\subsubsection{Phân tích so sánh định lượng}

Bảng~\ref{tab:codegen_comparison} tóm tắt so sánh ba phương pháp trên các tiêu chí quan trọng cho triển khai vi điều khiển bị giới hạn tài nguyên:

\begin{table}[H]
    \centering
    \caption{So sánh phương pháp triển khai mô hình ML trên vi điều khiển}
    \label{tab:codegen_comparison}
    \small
    \begin{tabular}{@{}lp{2.5cm}p{2.5cm}p{2.5cm}@{}}
        \toprule
        Tiêu chí & Runtime & Compiler & \textbf{Trực tiếp} \\
        \midrule
        Overhead flash & 100--300KB & 15--50KB & \textbf{0KB} \\
        Overhead suy luận & 15--40\% & $<$5\% & \textbf{0\%} \\
        Hỗ trợ RF/SVM & Hạn chế$^\dag$ & Không$^\ddag$ & \textbf{Đầy đủ} \\
        Minh bạch & Thấp & Trung bình & \textbf{Hoàn toàn} \\
        Cần ONNX/TFLite & Có & Có & \textbf{Không} \\
        Phức tạp triển khai & Thấp & Cao & Trung bình \\
        Hỗ trợ DNN sâu & Tốt & Tốt & Hạn chế \\
        \bottomrule
    \end{tabular}
    \vspace{4pt}
    \begin{minipage}{\linewidth}
    \footnotesize
    $^\dag$TFLite Micro chỉ hỗ trợ NN/CNN; RF/SVM không thể chuyển đổi trực tiếp do onnx2tf không hỗ trợ ONNX-ML TreeEnsembleClassifier. Keras surrogate có thể sử dụng nhưng là xấp xỉ (80--90\% tương đồng).\\
    $^\ddag$EON Compiler và TVM tập trung vào mạng nơ-ron; không hỗ trợ RF/SVM gốc.
    \end{minipage}
\end{table}

\subsubsection{Lập luận cho lựa chọn thiết kế}

Quyết định sử dụng tạo mã trực tiếp được lập luận bởi bốn yếu tố:

\paragraph{(1) Mức sử dụng bộ nhớ tối thiểu}
Trên các vi điều khiển lớp wearable như Seeed XIAO nRF52840 (1MB flash, 256KB RAM), mỗi KB đều quan trọng. TFLite Micro interpreter chiếm khoảng 150--250KB flash \cite{warden2019tinyml, ray2023tinyml_survey}---tương đương 15--25\% tổng flash khả dụng---chỉ riêng runtime, trước khi tính trọng số mô hình. Phương pháp trực tiếp loại bỏ hoàn toàn overhead này: một mô hình Neural Network chỉ chiếm 95KB (trọng số + logic suy luận + trích xuất đặc trưng), để lại $>$900KB cho firmware ứng dụng và giúp việc chuyển sang các board cùng lớp tài nguyên trở nên khả thi hơn.

\paragraph{(2) Minh bạch và khả năng xác minh}
Một yêu cầu quan trọng trong thiết kế hệ thống suy luận nhúng là \textit{khả năng kiểm tra và xác minh mã đã tạo} (xem Phần~\ref{sec:training_deployment_parity}). Khi phát hiện lỗi zero-padding và sai lệch công thức kurtosis, có thể truy vết trực tiếp đến mã C++ đã tạo, so sánh từng dòng với triển khai Python, và sửa chữa. Với runtime đóng gói (TFLite Micro) hoặc trình biên dịch đóng (EON Compiler), quá trình gỡ lỗi này sẽ không thể thực hiện---lỗi sẽ biểu hiện như ``mô hình hoạt động kém trên thiết bị'' mà không có cách xác định nguyên nhân gốc.

Khả năng xác minh này đặc biệt quan trọng cho ứng dụng y tế và an toàn, nơi các tiêu chuẩn quy định (FDA, IEC 62304) yêu cầu khả năng truy vết và kiểm toán đầy đủ cho mọi thành phần phần mềm liên quan đến quyết định lâm sàng.

\paragraph{(3) Hỗ trợ trực tiếp các thuật toán ML cổ điển}
Hệ sinh thái TFLite Micro và ONNX Runtime Micro được thiết kế chủ yếu cho mạng nơ-ron. Triển khai Random Forest hoặc SVM qua các nền tảng này yêu cầu chuyển đổi nhiều bước: scikit-learn $\rightarrow$ ONNX (\texttt{skl2onnx}) $\rightarrow$ TFLite (\texttt{onnx-tf}) $\rightarrow$ TFLite Micro. Mỗi bước chuyển đổi giới thiệu rủi ro sai lệch số và mất thông tin (ví dụ: ONNX không hỗ trợ chính xác tham số \texttt{class\_weight} của scikit-learn). Phương pháp trực tiếp đọc các thuộc tính nội bộ của scikit-learn (\texttt{model.coefs\_}, \texttt{model.support\_vectors\_}, \texttt{model.estimators\_}) mà không qua định dạng trung gian, loại bỏ hoàn toàn chuỗi chuyển đổi dễ lỗi này.

\paragraph{(4) Các thuật toán cổ điển là lựa chọn phù hợp cho HAR biên}
Lựa chọn tạo mã trực tiếp liên kết chặt chẽ với lựa chọn thuật toán. Cho bài toán HAR 5--10 lớp với $<$1000 vectors đặc trưng và 63 đặc trưng bất biến hướng, các thuật toán cổ điển đạt hiệu suất cạnh tranh (97--99\%) với chi phí tài nguyên thấp hơn đáng kể. Các mô hình này có cấu trúc suy luận đơn giản---nhân ma trận cho NN, duyệt cây cho RF, tính kernel cho SVM---có thể triển khai chính xác trong vài trăm dòng C++ mà không cần framework tính toán phức tạp. Điều này tạo thành một \textit{sweet spot} giữa hiệu suất mô hình và tính khả thi triển khai cho các thiết bị đeo bị giới hạn tài nguyên.

\subsubsection{Hạn chế và các phương pháp bổ sung}
Tạo mã trực tiếp không phải giải pháp phổ quát. Nó không phù hợp cho:
\begin{itemize}
    \item \textbf{Mạng sâu nhiều lớp}: LSTM, Transformer, hoặc CNN với hơn 4--5 lớp convolution yêu cầu các toán tử phức tạp (attention, bidirectional) mà triển khai thủ công dễ lỗi và khó bảo trì
    \item \textbf{Lượng tử hóa nâng cao}: INT8 inference với hiệu chỉnh phạm vi per-channel yêu cầu hạ tầng runtime mà TFLite Micro cung cấp sẵn
    \item \textbf{Mô hình thay đổi thường xuyên}: Nếu mô hình cần cập nhật OTA (over-the-air) mà không biên dịch lại firmware, runtime interpreter là lựa chọn tốt hơn
\end{itemize}
Đường dẫn TFLite Micro (Mục~\ref{sec:tflite_micro_deployment}) tồn tại song song với tạo mã trực tiếp; người dùng chọn phương pháp phù hợp: tạo mã trực tiếp cho ML cổ điển và CNN nhỏ, TFLite Micro cho các mô hình cần lượng tử hóa INT8 đầy đủ.

\subsection{Triển khai TFLite Micro}
\label{sec:tflite_micro_deployment}

Đường dẫn triển khai TFLite Micro là phương pháp bổ sung phù hợp cho các kiến trúc mô hình neural network. Quá trình chuyển đổi tuân theo pipeline ONNX $\to$ TensorFlow $\to$ TFLite với các hạn chế quan trọng về phạm vi hỗ trợ \cite{tensorflow2015_whitepaper}.

\subsubsection{Phạm vi hỗ trợ theo loại mô hình}
Nghiên cứu thực nghiệm cho thấy khả năng chuyển đổi khác nhau đáng kể giữa các loại mô hình:

\paragraph{Neural Network và CNN: Chuyển đổi trực tiếp (hỗ trợ)}
Các mô hình neural network (MLP) và CNN từ PyTorch có thể chuyển đổi thành công qua pipeline:
\begin{enumerate}
    \item PyTorch model $\rightarrow$ ONNX (\texttt{torch.onnx.export})
    \item ONNX $\rightarrow$ TensorFlow (\texttt{onnx-tf})
    \item TensorFlow $\rightarrow$ TFLite (\texttt{tf.lite.TFLiteConverter})
    \item TFLite $\rightarrow$ TFLite Micro (quantization INT8)
\end{enumerate}

\paragraph{Random Forest và SVM: Hạn chế ONNX-ML (không hỗ trợ trực tiếp)}
Các mô hình Random Forest và SVM gặp rào cản chuyển đổi do hạn chế trong hệ sinh thái ONNX-ML:
\begin{itemize}
    \item \textbf{Scikit-learn $\rightarrow$ ONNX}: \texttt{skl2onnx} tạo các toán tử ONNX-ML như \texttt{TreeEnsembleClassifier} và \texttt{SVMClassifier}
    \item \textbf{ONNX-ML $\rightarrow$ TensorFlow}: \texttt{onnx2tf} \textbf{không hỗ trợ} các toán tử ONNX-ML, chỉ hỗ trợ các toán tử ONNX chuẩn (operator domain "ai.onnx")
    \item \textbf{Kết quả}: Pipeline bị gián đoạn, không thể chuyển đổi RF/SVM sang TFLite
\end{itemize}

\subsubsection{Giải pháp Keras Surrogate cho RF/SVM}
Để khắc phục hạn chế ONNX-ML, \textit{phương pháp Keras surrogate} có thể được áp dụng:

\begin{enumerate}
    \item \textbf{Sinh dữ liệu huấn luyện surrogate}: Sử dụng mô hình RF/SVM gốc để tạo soft predictions trên tập dữ liệu đại diện
    \item \textbf{Huấn luyện Keras MLP}: Mạng neural đa lớp học ánh xạ từ features $\rightarrow$ soft predictions của RF/SVM
    \item \textbf{Chuyển đổi MLP $\rightarrow$ TFLite}: Keras MLP có thể chuyển đổi thành công qua pipeline tiêu chuẩn
\end{enumerate}

\textbf{Hạn chế Surrogate}: Keras surrogate là \textit{xấp xỉ}, không phải triển khai chính xác. Độ thỏa thuận kỳ vọng 80-90\% giữa mô hình gốc và surrogate. Cho độ chính xác tối đa, \textit{tạo mã trực tiếp C++ vẫn là phương pháp được khuyến nghị cho RF/SVM}.

\subsubsection{Khuyến nghị sử dụng}
\begin{itemize}
    \item \textbf{NN/CNN}: Sử dụng TFLite Micro để tận dụng runtime tối ưu hóa và hỗ trợ quantization INT8
    \item \textbf{RF/SVM}: Ưu tiên tạo mã trực tiếp C++ (chính xác, nhanh, nhỏ gọn). TFLite surrogate chỉ khi yêu cầu tích hợp runtime chuẩn hóa
\end{itemize}

\subsection{Cấu trúc mã đã tạo}
Đối với mỗi mô hình đã huấn luyện, trình tạo sản xuất một gói triển khai gồm các tệp tham số, tệp cài đặt và tệp ví dụ tích hợp. Cụ thể:
\begin{enumerate}
    \item \textbf{Tệp Header (.h)}: Chứa các tham số mô hình (ví dụ: vectors hỗ trợ, cấu trúc cây, ma trận trọng số), khai báo hàm, các macro số lượng đặc trưng/lớp
    \item \textbf{Tệp Triển khai (.c/.cpp/.py)}: Triển khai trích xuất đặc trưng, logic dự đoán (đặc thù cho thuật toán) và các hàm tiện ích, với phần mở rộng phụ thuộc nền tảng đích
    \item \textbf{Tệp Ví dụ Tích hợp}: Có thể là Arduino sketch (.ino), chương trình C/C++ ví dụ hoặc script MicroPython, minh họa khởi tạo cảm biến, quản lý bộ đệm cửa sổ trượt, vòng lặp suy luận thời gian thực và đầu ra dự đoán
\end{enumerate}

\textbf{Xác minh Công thức Thống kê}: Mã C++ đã tạo cho trích xuất đặc trưng sử dụng các công thức thống kê được xác minh khớp chính xác với triển khai Python (pandas/NumPy). Cụ thể, các đặc trưng kurtosis và skewness sử dụng \textit{độ lệch chuẩn tổng thể} (population standard deviation, chia $n$) cho chuẩn hóa z-scores, thay vì độ lệch chuẩn mẫu (chia $n-1$). Sự khác biệt này, mặc dù nhỏ ($\sim$1,3\% cho $n=150$), là hệ thống và tích lũy qua nhiều đặc trưng, có thể ảnh hưởng đến độ chính xác triển khai nếu không được xử lý.

\subsection{Lượng tử hóa mô hình cho triển khai biên}
\label{sec:quantization}

Lượng tử hóa (model quantization) là kỹ thuật nén mô hình then chốt trong triển khai TinyML, chuyển đổi các tham số mô hình từ biểu diễn dấu phẩy động 32-bit (float32) sang số nguyên 8-bit (int8) hoặc 16-bit (int16). Trên vi điều khiển như Cortex-M4F với 256~KB RAM và 1~MB flash, lượng tử hóa có thể quyết định liệu một mô hình có thể triển khai được hay không: chuyển đổi từ float32 sang int8 giảm kích thước trọng số $4\times$, giảm băng thông bộ nhớ và tăng tốc suy luận 2--3$\times$ nhờ các tập lệnh SIMD integer của ARM \cite{lin2022mcunet}.

\subsubsection{Các loại lượng tử hóa}

\paragraph{Lượng tử hóa sau huấn luyện (Post-Training Quantization, PTQ)}
PTQ áp dụng sau khi huấn luyện hoàn tất mà không cần sửa đổi quy trình huấn luyện. Mỗi tham số float32 $x$ được ánh xạ sang giá trị lượng tử hóa $x_q$ theo:
\begin{equation}
    x_q = \mathrm{clip}\!\left(\mathrm{round}\!\left(\frac{x}{s}\right) + z,\; 0,\; 255\right)
    \label{eq:quantization}
\end{equation}
trong đó $s = \frac{x_{\max} - x_{\min}}{255}$ là hệ số tỷ lệ và $z \in [0, 255]$ là zero-point giữ nguyên giá trị số không. Có hai biến thể PTQ: \textit{lượng tử hóa trọng số} (weight-only), trong đó chỉ trọng số mô hình được chuyển đổi còn phép tính suy luận vẫn dùng float32; và \textit{lượng tử hóa tĩnh} (static PTQ), trong đó cả trọng số lẫn giá trị kích hoạt đều được lượng tử hóa, đòi hỏi một tập hiệu chỉnh nhỏ (100--500 mẫu đại diện) để đo phân phối kích hoạt trong thực tế. PTQ thường gây suy giảm độ chính xác 0,5--2\% cho các mô hình HAR \cite{warden2019tinyml}.

\paragraph{Lượng tử hóa trong huấn luyện (Quantization-Aware Training, QAT)}
QAT chèn các toán tử \textit{fake quantization} (lượng tử hóa giả) vào đồ thị tính toán trong quá trình huấn luyện, mô phỏng sai số lượng tử hóa để mô hình học cách bù đắp. QAT thường giảm suy giảm độ chính xác xuống dưới 0,3\% nhưng đòi hỏi sửa đổi pipeline huấn luyện và thực tế chỉ áp dụng cho mạng nơ-ron---Random Forest và SVM có cấu trúc rời rạc (ngưỡng quyết định, vector hỗ trợ) không tương thích với kỹ thuật đạo hàm giả (straight-through estimator) mà QAT dựa vào.

\subsubsection{Chiến lược lượng tử hóa theo phương pháp tạo mã}

Tùy theo phương pháp tạo mã được chọn, lượng tử hóa được áp dụng theo hai đường dẫn khác nhau:

\paragraph{Lượng tử hóa trọng số trong tạo mã trực tiếp}
Trong đường dẫn tạo mã trực tiếp, lượng tử hóa được thực hiện ở cấp độ biểu diễn tham số trong mã C++ đã tạo. Chế độ tối ưu hóa kiểm soát số chữ số thập phân khi nhúng hằng số float32 vào mã nguồn---tương đương lượng tử hóa trọng số mức thấp:
\begin{itemize}
    \item \textbf{Accuracy}: float32 đầy đủ (6 chữ số thập phân), không xấp xỉ
    \item \textbf{Balanced}: 3 chữ số thập phân, tiết kiệm $\approx$30\% kích thước mã nguồn
    \item \textbf{Speed/Power}: 2 chữ số thập phân, tiết kiệm $\approx$45\% kích thước mã nguồn
\end{itemize}
Phép suy luận vẫn thực hiện bằng số học float32 trên FPU của Cortex-M4F---đây là lượng tử hóa \textit{lưu trữ}, không phải lượng tử hóa \textit{tính toán} đầy đủ. Đối với mạng nơ-ron, chế độ \textit{Power} có thể tùy chọn biểu diễn ma trận trọng số ở int16, thực hiện nhân tích lũy integer và chỉ dequantize một lần trước hàm kích hoạt:
\begin{equation}
    y_j = f\!\left(s \cdot \sum_i w_{ij}^{\mathrm{int16}} \cdot x_i^{\mathrm{int16}} + b_j\right)
\end{equation}
trong đó $s$ là tích của hai scale tương ứng (trọng số và đầu vào), $f(\cdot)$ là hàm kích hoạt ReLU hoặc softmax.

\paragraph{Lượng tử hóa INT8 đầy đủ qua đường dẫn TFLite Micro}
Đường dẫn TFLite Micro (áp dụng cho NN và CNN) hỗ trợ static PTQ với INT8 đầy đủ---cả trọng số lẫn kích hoạt. Quy trình thực hiện:
\begin{enumerate}
    \item Chuyển đổi mô hình Keras/PyTorch sang TFLite (\texttt{TFLiteConverter})
    \item Một tập hiệu chỉnh gồm 100--500 vectors đặc trưng đại diện được cung cấp để TFLiteConverter đo phân phối kích hoạt từng lớp
    \item Converter tự động tính $s$ và $z$ cho mỗi tensor theo Phương trình~\ref{eq:quantization}, tối ưu cho phạm vi kích hoạt thực tế
    \item Mô hình INT8 kết quả kết hợp với ARM CMSIS-NN \cite{banbury2021mlperf} để khai thác tập lệnh SIMD của Cortex-M4, thực thi 4 phép nhân int8 song song mỗi chu kỳ---tăng tốc 2--3$\times$ so với mô hình float32 tương đương
\end{enumerate}

Bảng~\ref{tab:quantization_tradeoff} so sánh hai đường dẫn lượng tử hóa theo các tiêu chí triển khai:

\begin{table}[H]
    \centering
    \caption{So sánh chiến lược lượng tử hóa theo phương pháp tạo mã}
    \label{tab:quantization_tradeoff}
    \small
    \begin{tabular}{@{}p{3.5cm}p{3.5cm}p{4.2cm}@{}}
        \toprule
        Tiêu chí & Tạo mã trực tiếp (Weight PTQ) & TFLite Micro (Static INT8 PTQ) \\
        \midrule
        Áp dụng cho & RF, SVM, NN, CNN & NN, CNN (không hỗ trợ RF/SVM) \\
        Giảm kích thước & 30--45\% (so với float32 đầy đủ) & $4\times$ (so với float32) \\
        Tăng tốc suy luận & Không đáng kể (vẫn dùng FPU float32) & 2--3$\times$ (SIMD int8) \\
        Suy giảm độ chính xác & $<$0,5\% & 0,5--2\% \\
        Yêu cầu tập hiệu chỉnh & Không & 100--500 mẫu \\
        Tính minh bạch & Hoàn toàn (mã có thể kiểm tra) & Hạn chế (runtime hộp đen) \\
        \bottomrule
    \end{tabular}
\end{table}

\subsubsection{Đặc thù lượng tử hóa cho đặc trưng HAR miền thời gian}

Bài toán HAR với đặc trưng thống kê miền thời gian có các đặc tính thuận lợi cho lượng tử hóa so với thị giác máy tính hay xử lý ngôn ngữ tự nhiên:

\begin{itemize}
    \item \textbf{Phạm vi đặc trưng có giới hạn vật lý}: Gia tốc IMU bị giới hạn ở $\pm$20~m/s², vận tốc góc ở $\pm$2000$^\circ$/s. Do đó, các đặc trưng thống kê (trung bình, RMS, độ lệch chuẩn) có phạm vi xác định, giúp xác định scale lượng tử hóa chính xác mà không cần tập hiệu chỉnh lớn.
    \item \textbf{Bền vững với nhiễu lượng tử hóa}: Random Forest quyết định theo ngưỡng nhị phân tại mỗi nút cây; sai số lượng tử hóa nhỏ hiếm khi đủ để đổi chiều quyết định phân nhánh, đặc biệt khi đặc trưng nằm xa ngưỡng. SVM nhạy hơn ở vùng ranh giới quyết định nhưng vẫn ổn định với int16.
    \item \textbf{Tham số scaler phải đồng bộ độ chính xác}: StandardScaler (trung bình $\mu$ và độ lệch chuẩn $\sigma$) được áp dụng trước mô hình. Để duy trì parity huấn luyện-triển khai, tham số scaler cần được nhúng với cùng độ chính xác với tham số mô hình---lượng tử hóa mô hình nhưng giữ scaler ở float32 đầy đủ sẽ gây ra sự không nhất quán hệ thống ảnh hưởng đến toàn bộ vector đặc trưng đầu vào.
    \item \textbf{Kurtosis và Skewness nhạy hơn}: Các đặc trưng hình dạng phân phối có phạm vi giá trị rộng hơn và nhạy với lượng tử hóa mạnh hơn các đặc trưng năng lượng (RMS, Energy). Do đó, chế độ \textit{Balanced} khuyến nghị sử dụng 3 chữ số thập phân để bảo toàn các đặc trưng này.
\end{itemize}

\subsection{Kiến trúc tạo mã đa kiến trúc}
\label{sec:multi_architecture_codegen}

Hai loại pipeline mô hình tồn tại với yêu cầu khác biệt:

\paragraph{Mô hình dựa trên đặc trưng (RF, SVM, NN)}
Nhận vector đặc trưng thống kê (33--90 chiều), áp dụng StandardScaler trên features đã trích xuất, sinh mã trích xuất đặc trưng đầy đủ trong C++.

\paragraph{Mô hình CNN end-to-end}
Nhận trực tiếp cửa sổ cảm biến thô ($150 \times 6 = 900$ đầu vào), chỉ áp dụng lọc tín hiệu (không trích xuất đặc trưng), StandardScaler trên dữ liệu thô.

Sự khác biệt này đòi hỏi kiến trúc tạo mã nhận biết loại mô hình (architecture-aware): metadata, logic xác minh và cấu trúc mã sinh ra đều phụ thuộc vào loại pipeline. Generator tự động phân biệt để sinh đúng cấu trúc cho từng loại.

\subsection{Các chiến lược tối ưu hóa}
Bốn chế độ tối ưu hóa được cấu hình khi tạo mã: \textit{Accuracy} (float32 đầy đủ, không đơn giản hóa); \textit{Speed} (giảm phức tạp mô hình, tính toán cache-friendly, mở rộng vòng lặp); \textit{Power} (số học điểm cố định int16, 2 chữ số thập phân, tận dụng ARM CMSIS-DSP); \textit{Balanced} (3 chữ số thập phân, độ chính xác hỗn hợp cho các tính toán khác nhau). Chi tiết so sánh tại Bảng~\ref{tab:optimization_modes}, Chương~\ref{chap:design_implementation}.

\subsection{Điều chỉnh đặc thù nền tảng}
Mã sinh ra được điều chỉnh theo họ thiết bị đích: Arduino-compatible (sinh .ino/.cpp/.h với Serial và Wire); ARM Cortex-M native (mã C độc lập, bộ nhớ tĩnh, tích hợp CMSIS); SDK/RTOS (include và entry-point cho ESP-IDF, Zephyr); MicroPython (module Python độc lập duy trì tương đồng công thức với C++). Nhờ phân lớp này, cùng một mô hình có thể xuất cho nhiều đích mà không cần huấn luyện lại.

Chi tiết triển khai phần mềm và phần cứng được trình bày trong Chương~\ref{chap:design_implementation}. Thiết lập thực nghiệm, tập dữ liệu và số liệu đánh giá được mô tả trong Chương~\ref{chap:results}.
